\begin{proof}[Proof of \cref{obj:thm:stable-grid-upper}]
Fix an experiment in \(\mathcal M_T^{(2)}(t_0,\zeta)\). All expectations below are under its observed trajectory law, denoted by \(\mathbb E_M\). Use the moments of \cref{obj:def:observable-moments} and the joint-state quantities of \cref{obj:lem:observable-intervention-moments}.

\begin{enumerate}
\item Use zero-based epochs throughout this proof: epoch \(t\) here corresponds to epoch \(t+1\) in \cref{obj:def:observable-moments}. Accordingly, define the shifted scores by
\[
Z_t(k):=Y_{t+1}\prod_{j=t-k+1}^{t+1}\rho_j,
\qquad 0\le k\le6,\quad k\le t<T,
\]
where the reward and likelihood ratios on the right retain their one-based indexing. Thus \(Z_t(k)\) is the score at one-based epoch \(t+1\), agreeing with the zero-based convention of \cref{obj:lem:observable-intervention-moments}.

Define the retained epoch set, its cardinality, and the maximum empirical moment error by
\[
\mathcal I_T:=\{6,\ldots,T-1\},
\qquad
n_T:=|\mathcal I_T|=T-6,
\qquad
\epsilon_T(w):=
\max_{0\le k\le6}|\widehat m_k(w)-m_k|,
\]
where \(w\in(\mathcal X\times\mathcal A\times\mathbb R)^T\) is an observed trajectory. Changing variables in the one-based retained-window sum gives
\[
\widehat m_k
=
\frac1{T-6}\sum_{t=6}^{T-1}Z_t(k)
=
\frac1{n_T}\sum_{t\in\mathcal I_T}Z_t(k),
\qquad 0\le k\le6.
\]
The parameter conditions give
\[
0<\alpha=\exp(-1/t_0)<1,
\qquad
L=\exp(\zeta)\ge1,
\qquad
B_\alpha\ge0,
\qquad
V_{\alpha,L}\ge0,
\]
and the horizon condition gives
\[
T>0,\qquad n_T>0,\qquad T\le2n_T,\qquad T\ge1.
\]
The maximum defining \(\epsilon_T(w)\) is attained among seven indices. Its square is therefore one of the nonnegative summands on the right-hand side of
\[
\epsilon_T(w)^2
\le
\sum_{k=0}^{6}(\widehat m_k(w)-m_k)^2.
\]
The observed law is a probability law, and the empirical moments are square-integrable. Thus the sum on the right is integrable; measurability of the finite maximum and the displayed domination also give integrability of \(\epsilon_T^2\).
% lean: CausalSmith.Stat.PomdpBinaryhiddenRate.seven_error_max_sq_le_sum
% lean: CausalSmith.Stat.PomdpBinaryhiddenRate.stableGrid_upper

\item For \(0\le k\le6\) and \(t,t'\in\mathcal I_T\), define
\[
c_{k;t,t'}
:=
\mathbb E_M\!\left[
  (Z_t(k)-m_k)(Z_{t'}(k)-m_k)
\right].
\]
Every retained epoch satisfies \(k\le t\). Consequently, \cref{obj:lem:observable-intervention-moments} supplies
\[
\mathbb E_M[Z_t(k)]=m_k=d_bP_e^kr_e,
\qquad
\mathbb E_M[Z_t(k)^2]\le L^{k+1}.
\]
In particular,
\[
0\le c_{k;t,t}
=
\mathbb E_M[(Z_t(k)-m_k)^2]
=
\mathbb E_M[Z_t(k)^2]-m_k^2
\le L^{k+1}.
\]
For either choice of sign, nonnegativity of the centered square gives
\[
\begin{aligned}
0
&\le
\mathbb E_M\!\left[
 \bigl((Z_t(k)-m_k)\pm(Z_{t'}(k)-m_k)\bigr)^2
\right]\\
&=
c_{k;t,t}+c_{k;t',t'}\pm2c_{k;t,t'}.
\end{aligned}
\]
Using both signs yields
\[
|c_{k;t,t'}|
\le
\frac{c_{k;t,t}+c_{k;t',t'}}2
\le L^{k+1}.
\]
For separated epochs, the covariance conclusion of
\cref{obj:lem:observable-intervention-moments} additionally gives
\[
|c_{k;t,t'}|
\le\alpha^{\,t'-t-k-1}
\qquad\text{when }t'-t\ge k+1.
\]
These identities and bounds apply under the observed law because the scores are functions of the observed trajectory.
% lean: CausalSmith.Stat.PomdpBinaryhiddenRate.observable_score_variance_le
% lean: CausalSmith.Stat.PomdpBinaryhiddenRate.observable_score_covariance_le
% lean: CausalSmith.Stat.PomdpBinaryhiddenRate.observable_score_covariance_le_separated

\item Fix \(0\le k\le6\) and \(t\in\mathcal I_T\). Partition its future retained epochs according to whether \(1\le t'-t\le k\) or \(t'-t\ge k+1\). There are at most \(k\) epochs in the first part. In the second part, the exponents \(t'-t-k-1\) are distinct integers in \(\{0,\ldots,T-1\}\). Hence
\[
\begin{aligned}
\sum_{\substack{t'\in\mathcal I_T\\t<t'}}
 |c_{k;t,t'}|
&=
\sum_{\substack{t'\in\mathcal I_T\\1\le t'-t\le k}}
 |c_{k;t,t'}|
+
\sum_{\substack{t'\in\mathcal I_T\\t'-t\ge k+1}}
 |c_{k;t,t'}|\\
&\le
kL^{k+1}+\sum_{j=0}^{T-1}\alpha^j\\
&=
kL^{k+1}+\frac{1-\alpha^T}{1-\alpha}\\
&\le
kL^{k+1}+\frac1{1-\alpha}.
\end{aligned}
\]
For \(k=0\), the first sum is empty and equals zero. Symmetry of covariance allows the full absolute covariance sum to be split into its diagonal and twice its upper triangle:
\[
\begin{aligned}
\sum_{t\in\mathcal I_T}\sum_{t'\in\mathcal I_T}c_{k;t,t'}
&\le
\sum_{t\in\mathcal I_T}\sum_{t'\in\mathcal I_T}|c_{k;t,t'}|\\
&=
\sum_{t\in\mathcal I_T}|c_{k;t,t}|
+
2\sum_{t\in\mathcal I_T}
 \sum_{\substack{t'\in\mathcal I_T\\t<t'}}
 |c_{k;t,t'}|\\
&\le
n_TL^{k+1}
+
2n_T\left(kL^{k+1}+\frac1{1-\alpha}\right)\\
&=
n_T\left((1+2k)L^{k+1}+\frac2{1-\alpha}\right)\\
&\le
n_T\left(13L^7+\frac2{1-\alpha}\right)
=
n_TV_{\alpha,L}.
\end{aligned}
\]
The last inequality uses \(k\le6\) and \(L\ge1\).

Averaging the score means over the common window gives
\[
\mathbb E_M[\widehat m_k]
=
\frac1{n_T}\sum_{t\in\mathcal I_T}\mathbb E_M[Z_t(k)]
=
m_k.
\]
Expanding the centered square of this finite average therefore yields
\[
\begin{aligned}
\mathbb E_M[(\widehat m_k-m_k)^2]
&=
\frac1{n_T^2}
\sum_{t\in\mathcal I_T}\sum_{t'\in\mathcal I_T}c_{k;t,t'}\\
&\le
\frac{V_{\alpha,L}}{n_T}.
\end{aligned}
\]
% lean: CausalSmith.Stat.PomdpBinaryhiddenRate.binary_sum_abs_covariance_future_le
% lean: CausalSmith.Stat.PomdpBinaryhiddenRate.binary_sum_abs_symmetric_eq_diag_add_two_upper
% lean: CausalSmith.Stat.PomdpBinaryhiddenRate.empiricalMoment_mse_eq_covariance_sum
% lean: CausalSmith.Stat.PomdpBinaryhiddenRate.empiricalMoment_mse_le_varianceFactor

\item Integrating \(\epsilon_T(w)^2\le\sum_{k=0}^{6}(\widehat m_k(w)-m_k)^2\) and applying the seven moment bounds gives
\[
\mathbb E_M[\epsilon_T^2]
\le
\sum_{k=0}^{6}\mathbb E_M[(\widehat m_k-m_k)^2]
\le
\frac{7V_{\alpha,L}}{n_T}.
\]
% lean: CausalSmith.Stat.PomdpBinaryhiddenRate.empiricalMoment_max_mse_le_varianceFactor

\item We next obtain a value bound for the selected grid realization. To distinguish the grid resolution from the fixed experiment, write
\[
M_{\mathrm{grid}}
:=
\left\lceil(22+36/\alpha)T\right\rceil,
\qquad
U:=\frac{\mathbf1\mathbf1^{\mathsf T}}4,
\qquad
\lambda:=\frac{\alpha}{\alpha+6/M_{\mathrm{grid}}}.
\]
Here \(\mathbf1\) is the four-dimensional vector of ones. Thus
\[
M_{\mathrm{grid}}>0,
\qquad
0<\lambda<1,
\qquad
P(R)=\lambda R+(1-\lambda)U.
\]
Use a fixed ordering of the four joint states as the four grid coordinates; this relabeling preserves scalar matrix moments and total variation.

The candidate set of \cref{obj:def:stable-grid-estimator} is finite and nonempty. Indeed,
\[
\bigl((1,0,0,0),\ \mathbf1(1,0,0,0),\ (0,0,0,0)\bigr)
\in\mathcal C_{M_{\mathrm{grid}}},
\]
because all its coordinates lie on the required grids and
\[
\delta\bigl(\mathbf1(1,0,0,0)\bigr)=0
\le\alpha+\frac6{M_{\mathrm{grid}}}.
\]
The lexicographic minimizer therefore belongs to this candidate set.

For every candidate \((\nu,R,r)\), its stabilized matrix is row-stochastic and satisfies
\[
P(R)_{ij}\ge\frac{1-\lambda}{4}.
\]
The common uniform component cancels between rows, so
\[
\begin{aligned}
\delta(P(R))
&=
\lambda\delta(R)\\
&\le
\lambda\left(\alpha+\frac6{M_{\mathrm{grid}}}\right)
=
\alpha.
\end{aligned}
\]
The same uniform minorization gives, for probability rows \(\mu,\mu'\),
\[
\|\mu P(R)-\mu'P(R)\|_{\mathrm{TV}}
\le
\lambda\|\mu-\mu'\|_{\mathrm{TV}}.
\]
Since \(\lambda<1\), contraction on the finite probability simplex supplies a stationary probability row:
\[
\pi(P(R))P(R)=\pi(P(R)),
\qquad
\pi(P(R))_i\ge0,
\qquad
\sum_{i=1}^{4}\pi(P(R))_i=1.
\]

Write the selected quantities as in \cref{obj:def:stable-grid-estimator}:
\[
(\nu_*,R_*,r_*)\in\mathcal C_{M_{\mathrm{grid}}},
\qquad
P_*:=P(R_*),
\qquad
\widehat\theta_T=\pi(P_*)r_*.
\]
The experiment's class conditions give a row-stochastic \(P_e\), probability rows \(d_b,d_e\), and a unit-range reward vector:
\[
P_e(i,j)\ge0,\quad \sum_jP_e(i,j)=1,
\qquad
d_b(i),d_e(i)\ge0,\quad \sum_id_b(i)=\sum_id_e(i)=1,
\]
\[
d_eP_e=d_e,
\qquad
0\le r_e(i)\le1.
\]
Applying \cref{obj:ass:target-contraction} to point masses gives
\[
\delta(P_e)
=
\max_{i,j}\frac12\sum_s|P_e(i,s)-P_e(j,s)|
\le\alpha.
\]
Together with \(\theta=d_er_e\) from \cref{obj:def:target-value}, these facts discharge the hypotheses of \cref{obj:lem:pair-polynomial-modulus}. Its value bound is
\[
\begin{aligned}
|\widehat\theta_T-\theta|
&\le
B_\alpha
\max_{0\le k\le6}
|\nu_*P_*^kr_*-d_bP_e^kr_e|\\
&=
B_\alpha
\max_{0\le k\le6}
|\nu_*P_*^kr_*-m_k|.
\end{aligned}
\]
% lean: CausalSmith.Stat.PomdpBinaryhiddenRate.selectGridCode_mem
% lean: CausalSmith.Stat.PomdpBinaryhiddenRate.stabilizedGridMatrix_stationary
% lean: CausalSmith.Stat.PomdpBinaryhiddenRate.stabilizedGridMatrix_dobrushin_le
% lean: CausalSmith.Stat.PomdpBinaryhiddenRate.grid_value_modulus

\item Construct an approximating candidate by rounding. For every probability row \(\mu\) on the four coordinates, define
\[
\bigl(\mathcal R_{\mathrm{grid}}(\mu)\bigr)_i
:=
\begin{cases}
\lfloor M_{\mathrm{grid}}\mu_i\rfloor/M_{\mathrm{grid}},
 &i=1,2,3,\\[2mm]
1-\displaystyle\sum_{j=1}^{3}
 \lfloor M_{\mathrm{grid}}\mu_j\rfloor/M_{\mathrm{grid}},
 &i=4.
\end{cases}
\]
The first three coordinates are rounded down, and their sum is at most one. The fourth coordinate is therefore nonnegative, is a grid multiple, and makes the total mass one. For \(i=1,2,3\),
\[
0\le
\mu_i-\bigl(\mathcal R_{\mathrm{grid}}(\mu)\bigr)_i
\le\frac1{M_{\mathrm{grid}}},
\]
while
\[
\mu_4-\bigl(\mathcal R_{\mathrm{grid}}(\mu)\bigr)_4
=
-\sum_{i=1}^{3}
 \left(\mu_i-\bigl(\mathcal R_{\mathrm{grid}}(\mu)\bigr)_i\right).
\]
Consequently,
\[
\begin{aligned}
\|\mu-\mathcal R_{\mathrm{grid}}(\mu)\|_{\mathrm{TV}}
&=
\frac12\sum_{i=1}^{4}
 \left|\mu_i-\bigl(\mathcal R_{\mathrm{grid}}(\mu)\bigr)_i\right|\\
&=
\sum_{i=1}^{3}
 \left(\mu_i-\bigl(\mathcal R_{\mathrm{grid}}(\mu)\bigr)_i\right)
\le\frac3{M_{\mathrm{grid}}}.
\end{aligned}
\]

Define the rounded initial row, transition rows, and rewards by
\[
\nu_{\mathrm{grid}}:=\mathcal R_{\mathrm{grid}}(d_b),
\qquad
R_{\mathrm{grid}}(i,\cdot)
:=
\mathcal R_{\mathrm{grid}}(P_e(i,\cdot)),
\qquad
r_{\mathrm{grid}}(i)
:=
\frac{\lfloor M_{\mathrm{grid}}r_e(i)\rfloor}{M_{\mathrm{grid}}}.
\]
They satisfy
\[
\|\nu_{\mathrm{grid}}-d_b\|_{\mathrm{TV}}
\le\frac3{M_{\mathrm{grid}}},
\qquad
\max_i\|R_{\mathrm{grid}}(i,\cdot)-P_e(i,\cdot)\|_{\mathrm{TV}}
\le\frac3{M_{\mathrm{grid}}},
\]
\[
0\le r_{\mathrm{grid}}(i)\le1,
\qquad
0\le r_e(i)-r_{\mathrm{grid}}(i)
\le\frac1{M_{\mathrm{grid}}}.
\]
For each pair of rows, the triangle inequality gives
\[
\begin{aligned}
\|R_{\mathrm{grid}}(i,\cdot)-R_{\mathrm{grid}}(j,\cdot)\|_{\mathrm{TV}}
&\le
\|R_{\mathrm{grid}}(i,\cdot)-P_e(i,\cdot)\|_{\mathrm{TV}}\\
&\quad+
\|P_e(i,\cdot)-P_e(j,\cdot)\|_{\mathrm{TV}}\\
&\quad+
\|P_e(j,\cdot)-R_{\mathrm{grid}}(j,\cdot)\|_{\mathrm{TV}}\\
&\le\alpha+\frac6{M_{\mathrm{grid}}}.
\end{aligned}
\]
Taking the maximum over the row pair proves
\[
\delta(R_{\mathrm{grid}})
\le\alpha+\frac6{M_{\mathrm{grid}}},
\qquad
(\nu_{\mathrm{grid}},R_{\mathrm{grid}},r_{\mathrm{grid}})
\in\mathcal C_{M_{\mathrm{grid}}}.
\]
% lean: CausalSmith.Stat.PomdpBinaryhiddenRate.reward_grid_rounding
% lean: CausalSmith.Stat.PomdpBinaryhiddenRate.simplex_grid_rounding
% lean: CausalSmith.Stat.PomdpBinaryhiddenRate.gridDobrushin_perturbation
% lean: CausalSmith.Stat.PomdpBinaryhiddenRate.grid_candidate_moment_approximation

\item Define the stabilized approximating matrix by
\[
P_{\mathrm{grid}}
:=
P(R_{\mathrm{grid}})
=
\lambda R_{\mathrm{grid}}+(1-\lambda)U.
\]
For every row, its distance from the uniform probability row is at most one. Also,
\[
1-\lambda
=
\frac{6/M_{\mathrm{grid}}}{\alpha+6/M_{\mathrm{grid}}}
\le\frac6{\alpha M_{\mathrm{grid}}}.
\]
It follows that
\[
\begin{aligned}
\|R_{\mathrm{grid}}(i,\cdot)-P_{\mathrm{grid}}(i,\cdot)\|_{\mathrm{TV}}
&=
(1-\lambda)
\|R_{\mathrm{grid}}(i,\cdot)-U(i,\cdot)\|_{\mathrm{TV}}\\
&\le\frac6{\alpha M_{\mathrm{grid}}},
\end{aligned}
\]
and hence
\[
\|P_e(i,\cdot)-P_{\mathrm{grid}}(i,\cdot)\|_{\mathrm{TV}}
\le
\frac{3+6/\alpha}{M_{\mathrm{grid}}}.
\]

For every integer \(k\ge0\), stochasticity and the reward bounds imply
\[
0\le(P_e^kr_e)(i)\le1.
\]
Integration of a \([0,1]\)-valued function against the difference of two probability rows is bounded by their total variation distance: centering that function at \(1/2\) gives this bound directly. We claim
\[
\|P_e^kr_e-P_{\mathrm{grid}}^kr_{\mathrm{grid}}\|_\infty
\le
\frac{k(3+6/\alpha)+1}{M_{\mathrm{grid}}},
\qquad k\ge0.
\]
For \(k=0\), this is the reward-rounding bound. At the induction step, insert
\(P_{\mathrm{grid}}P_e^kr_e\) between the two expressions. The row approximation bound controls the first difference, and averaging the previous error with a stochastic row controls the second:
\[
\begin{aligned}
&\|P_e^{k+1}r_e-P_{\mathrm{grid}}^{k+1}r_{\mathrm{grid}}\|_\infty\\
&\quad\le
\|(P_e-P_{\mathrm{grid}})P_e^kr_e\|_\infty
+
\|P_{\mathrm{grid}}
 (P_e^kr_e-P_{\mathrm{grid}}^kr_{\mathrm{grid}})\|_\infty\\
&\quad\le
\frac{3+6/\alpha}{M_{\mathrm{grid}}}
+
\|P_e^kr_e-P_{\mathrm{grid}}^kr_{\mathrm{grid}}\|_\infty\\
&\quad\le
\frac{(k+1)(3+6/\alpha)+1}{M_{\mathrm{grid}}}.
\end{aligned}
\]
The initial-distribution error can now be separated from the power-and-reward error:
\[
\begin{aligned}
|m_k-\nu_{\mathrm{grid}}P_{\mathrm{grid}}^kr_{\mathrm{grid}}|
&=
|d_bP_e^kr_e-\nu_{\mathrm{grid}}P_{\mathrm{grid}}^kr_{\mathrm{grid}}|\\
&\le
|(d_b-\nu_{\mathrm{grid}})P_e^kr_e|
+
|\nu_{\mathrm{grid}}
 (P_e^kr_e-P_{\mathrm{grid}}^kr_{\mathrm{grid}})|\\
&\le
\frac3{M_{\mathrm{grid}}}
+
\frac{k(3+6/\alpha)+1}{M_{\mathrm{grid}}}\\
&=
\frac{4+k(3+6/\alpha)}{M_{\mathrm{grid}}}.
\end{aligned}
\]
For \(0\le k\le6\), the coefficient and ceiling calibration give
\[
4+k(3+6/\alpha)\le22+36/\alpha,
\qquad
(22+36/\alpha)T\le M_{\mathrm{grid}},
\]
so
\[
|m_k-\nu_{\mathrm{grid}}P_{\mathrm{grid}}^kr_{\mathrm{grid}}|
\le
\frac{22+36/\alpha}{M_{\mathrm{grid}}}
\le\frac1T.
\]
% lean: CausalSmith.Stat.PomdpBinaryhiddenRate.stabilizedGridMatrix_row_error
% lean: Causalean.Mathlib.Probability.FiniteMarkovPerturbation.stochastic_power_reward_error
% lean: Causalean.Mathlib.Probability.FiniteMarkovPerturbation.stochastic_moment_error
% lean: CausalSmith.Stat.PomdpBinaryhiddenRate.grid_approximation

\item Fix an observed trajectory \(w\). The approximating candidate has empirical discrepancy bounded by
\[
\begin{aligned}
\max_{0\le k\le6}
|\widehat m_k(w)-\nu_{\mathrm{grid}}P_{\mathrm{grid}}^kr_{\mathrm{grid}}|
&\le
\max_{0\le k\le6}|\widehat m_k(w)-m_k|\\
&\quad+
\max_{0\le k\le6}
|m_k-\nu_{\mathrm{grid}}P_{\mathrm{grid}}^kr_{\mathrm{grid}}|\\
&\le\epsilon_T(w)+T^{-1}.
\end{aligned}
\]
Because the selected candidate minimizes this discrepancy,
\[
\max_{0\le k\le6}
|\widehat m_k(w)-\nu_*P_*^kr_*|
\le\epsilon_T(w)+T^{-1}.
\]
Applying the triangle inequality once more gives
\[
\begin{aligned}
\max_{0\le k\le6}|\nu_*P_*^kr_*-m_k|
&\le
\max_{0\le k\le6}|\nu_*P_*^kr_*-\widehat m_k(w)|
+\epsilon_T(w)\\
&\le2\epsilon_T(w)+T^{-1}.
\end{aligned}
\]
Combining this with the value modulus obtained above yields
\[
|\widehat\theta_T(w)-\theta|
\le
B_\alpha(2\epsilon_T(w)+T^{-1}).
\]
The right-hand side is nonnegative, so squaring preserves the inequality. Moreover,
\((2\epsilon_T(w)-T^{-1})^2\ge0\) gives
\[
\begin{aligned}
(\widehat\theta_T(w)-\theta)^2
&\le
B_\alpha^2(2\epsilon_T(w)+T^{-1})^2\\
&\le
B_\alpha^2\left(8\epsilon_T(w)^2+2T^{-2}\right).
\end{aligned}
\]
% lean: CausalSmith.Stat.PomdpBinaryhiddenRate.selected_grid_moment_error
% lean: CausalSmith.Stat.PomdpBinaryhiddenRate.selected_grid_intervention_moment_error
% lean: CausalSmith.Stat.PomdpBinaryhiddenRate.stableGridEstimator_value_error
% lean: CausalSmith.Stat.PomdpBinaryhiddenRate.stableGrid_upper

\item The sampling and grid terms admit separate horizon bounds. Since
\(V_{\alpha,L}\ge0\) and \(T\le2n_T\),
\[
56V_{\alpha,L}T
\le112V_{\alpha,L}n_T.
\]
Dividing by the positive product \(n_TT\) gives
\[
8\left(\frac{7V_{\alpha,L}}{n_T}\right)
=
\frac{56V_{\alpha,L}}{n_T}
\le
\frac{112V_{\alpha,L}}T.
\]
Also, \(T\ge1\) implies \(T^2\ge T>0\), and therefore
\[
T^{-2}\le T^{-1}.
\]
% lean: CausalSmith.Stat.PomdpBinaryhiddenRate.stableGrid_upper

\item The pointwise upper bound is integrable by integrability of \(\epsilon_T^2\). Monotonicity of integration for the nonnegative squared loss, followed by linearity of expectation under a probability law, now gives
\[
\begin{aligned}
\mathbb E_M[(\widehat\theta_T-\theta)^2]
&\le
\mathbb E_M\!\left[
 B_\alpha^2(8\epsilon_T^2+2T^{-2})
\right]\\
&=
B_\alpha^2\left(8\mathbb E_M[\epsilon_T^2]+2T^{-2}\right)\\
&\le
B_\alpha^2\left(
 8\frac{7V_{\alpha,L}}{n_T}+2T^{-2}
\right)\\
&\le
B_\alpha^2\left(
 \frac{112V_{\alpha,L}}T+\frac2T
\right)\\
&=
\frac{B_\alpha^2(112V_{\alpha,L}+2)}T.
\end{aligned}
\]
The fixed experiment was arbitrary, and every bound uses the same constants throughout the stated class.
% lean: CausalSmith.Stat.PomdpBinaryhiddenRate.stableGrid_upper
\end{enumerate}
\end{proof}
